\documentclass[9pt,a4paper,twocolumn,twoside]{tau-class/tau}
\usepackage[spanish,es-nodecimaldot,es-noindentfirst]{babel}
\usepackage{graphicx}
\usepackage{url}

\doctype{Reporte Ejecutivo Integrado}
\title{Reporte de performance y panorama competitivo: Psicotelcon}

\author[a,1]{Equipo de Analitica Oderbiz}
\affil[a]{Unidad de Inteligencia Comercial}

\dates{Compilado el \today}
\docinfo{Documento generado con datos internos \texttt{llm\_context\_report.page.v1} y evidencia publica verificable. Toda inferencia se etiqueta como \texttt{inferencia}.}
\footinfo{Reporte Oderbiz}
\theday{\today}
\organization{Oderbiz}
\leadauthor{Equipo Oderbiz}

\begin{abstract}
En la ventana \texttt{last\_30d}, la cuenta de Psicotelcon registra \$21.30 de inversion, 14,759 impresiones, 12,695 de alcance y CTR de 1.57\% (\texttt{observado}). El embudo muestra 209 clics unicos, 11 conversaciones iniciadas y 11 primeras respuestas, con costo por conversacion estimado de \$1.94 (\texttt{estimado}). La concentracion geografica en Loja absorbe la mayor parte del presupuesto y volumen de resultados de mensajeria (\texttt{observado}). En competencia local, se observa mezcla de centros especializados, profesionales independientes y directorios de comparacion por precio/servicio (\texttt{observado}). La prioridad estrategica es mejorar trazabilidad de conversion y explotar creatividades orientadas a motivos de consulta de alta intencion (\texttt{inferencia}).
\end{abstract}

\keywords{Meta Ads, psicologia infantil, Loja, performance marketing, inteligencia competitiva}

\begin{document}
\maketitle
\thispagestyle{firststyle}

\section{Resumen ejecutivo}
\begin{itemize}
  \item Performance eficiente en awareness/interaccion con CPM de \$1.44 y CTR de 1.57\% (\texttt{observado}).
  \item Embudo conversacional funcional: 11 conversaciones y 11 primeras respuestas sobre 209 clics unicos (5.26\% clic unico a conversacion, \texttt{estimado}).
  \item Alto sesgo geografico hacia Loja: \$20.53 de \$21.30 del gasto total (96.4\%, \texttt{estimado}).
  \item Tramo etario 25--34 concentra mayor volumen de conexiones de mensajeria con mejor eficiencia relativa en costo por resultado (\texttt{observado}).
  \item Riesgo principal: baja salida a propiedades externas (\texttt{outbound\_clicks}=0), limitando atribucion fuera de Meta (\texttt{observado}+\texttt{inferencia}).
\end{itemize}

\section{Contexto del negocio y periodo analizado}
\textbf{Pagina}: Centro de Psicologia Integral y Terapia - Psicotelcon (\texttt{page\_id}: 1506380769434870).\\
\textbf{Cuenta}: Dra Marizol Jimenez Luzon (\texttt{account\_id}: act\_131112367482947).\\
\textbf{Periodo}: \texttt{date\_preset=last\_30d}, \texttt{date\_range.since=null}, \texttt{date\_range.until=null}.\\
\textbf{Campana filtrada}: no aplica (\texttt{filters.campaign\_id=null}, \texttt{filters.campaign\_name=null}).\\
\textbf{Moneda}: USD.\\
\textbf{Zona horaria}: America/Guayaquil.

El inventario de campanas disponibles confirma dominancia de objetivos \texttt{MESSAGES}, con apoyo de \texttt{OUTCOME\_ENGAGEMENT} y casos puntuales de trafico (\texttt{observado}).

\section{Baseline interno}
\subsection{KPIs base}
\begin{itemize}
  \item Spend: \$21.30; impresiones: 14,759; reach: 12,695; CPM: \$1.44; CTR: 1.57\% (\texttt{observado}).
  \item Funnel: 209 clics unicos, 11 conversaciones iniciadas, 11 primeras respuestas (\texttt{observado}).
  \item Costo por clic unico: \$0.10; costo por conversacion iniciada: \$1.94 (\texttt{estimado}).
  \item Calidad de trafico: \texttt{outbound\_clicks}=0 y \texttt{cost\_per\_outbound\_click}=0 por ausencia de eventos (\texttt{observado}).
\end{itemize}

\subsection{Tendencia temporal (dias con entrega)}
La serie diaria presenta 6 dias con gasto efectivo. El mejor CTR diario fue 1.88\% (2026-04-20) y el menor 0.98\% (2026-04-17), con CPM en rango \$1.19--\$1.67 (\texttt{observado}). La entrega parece concentrada en pulsos de activacion y no en pauta uniforme continua (\texttt{inferencia}).

\subsection{Embudo y calidad}
Tasa de clic unico a conversacion: 11/209 = 5.26\% (\texttt{estimado}). La paridad entre conversaciones iniciadas y primeras respuestas sugiere buena respuesta operativa inicial del canal de mensajeria (\texttt{inferencia-soporte}). Sin embargo, la ausencia de eventos outbound limita lectura de intencion a etapas mas profundas fuera de la plataforma (\texttt{observado}+\texttt{inferencia}).

\subsection{Sena geo y demografica}
\textbf{Geo}: Loja concentra 10 resultados de mensajeria con CPA \$2.05; Zamora Chinchipe aporta 1 resultado con CPA \$0.77 (\texttt{observado}).

\textbf{Edad}: 25--34 concentra 7 conexiones de mensajeria (costo por resultado \$0.98), mientras 35--44 y 45--54 muestran costo por resultado significativamente mayor en este corte (\texttt{observado}).

\section{Panorama competitivo}
\subsection{Senales competitivas observadas}
\begin{itemize}
  \item Directorios y rankings de psicologos en Loja incrementan comparacion abierta entre especialistas y centros (\texttt{observado}).
  \item Existen propuestas locales con foco en psicopedagogia, lenguaje y neurodesarrollo en formato centro especializado (\texttt{observado}).
  \item Se observan referencias publicas de precios por sesion/paquetes en algunos actores del mercado, lo que puede anclar expectativas de valor (\texttt{observado}).
\end{itemize}

\subsection{Lectura estrategica}
La competencia no se limita a clinicas directas; tambien compiten marketplaces y perfiles individuales bien indexados. Diferenciar por especializacion infantil + proceso terapeutico + evidencia de resultados es clave para defender conversion en audiencias de padres (\texttt{inferencia}).

\section{Diagnostico integrado (fortalezas, brechas, riesgos, oportunidades)}
\subsection{Fortalezas}
\begin{itemize}
  \item Eficiencia media-alta de compra de medios para volumen local (CPM bajo, CTR saludable) (\texttt{observado}).
  \item Buen comportamiento de mensajeria inicial (11/11 de inicio a primera respuesta en el corte) (\texttt{observado}).
  \item Posicionamiento en nicho infantil/psicopedagogico con huella publica verificable (\texttt{observado}).
\end{itemize}

\subsection{Brechas}
\begin{itemize}
  \item Sin trazabilidad de conversion downstream (cita confirmada, asistencia, continuidad) en este dataset (\texttt{observado}).
  \item Nomenclatura de campanas heterogenea, con varios IDs sin naming estrategico (\texttt{observado}).
  \item Baja senal de trafico a activos propios fuera de Meta en el periodo (\texttt{observado}).
\end{itemize}

\subsection{Riesgos}
\begin{itemize}
  \item Dependencia excesiva de mensajeria in-platform para medir valor real del lead (\texttt{inferencia}).
  \item Entorno competitivo con comparacion publica por precio/listados que puede comprimir conversion por diferenciacion insuficiente (\texttt{inferencia-soporte}).
\end{itemize}

\subsection{Oportunidades}
\begin{itemize}
  \item Escalar 25--34 con creatividades por problema especifico: lenguaje, conducta, aprendizaje (\texttt{inferencia}).
  \item Diseñar embudo mixto mensajes + destino trackeable para atribucion de calidad (\texttt{inferencia}).
  \item Replicar cobertura eficiente de Loja hacia provincias colindantes con pruebas controladas (\texttt{inferencia}).
\end{itemize}

\section{Plan de accion priorizado (P1/P2/P3)}
\subsection{P1: alto impacto / bajo esfuerzo}
\begin{itemize}
  \item Estandarizar naming de campanas/adsets/ads: objetivo--audiencia--servicio--fecha.
  \item Implementar tablero semanal con 4 KPI: costo por conversacion iniciada, costo por primera respuesta, tasa clic unico a conversacion, tiempo de primera respuesta.
  \item Lanzar 3 variaciones de copy por dolor principal (lenguaje, conducta, aprendizaje) sobre los anuncios activos.
  \item Definir script de primer contacto para mover de chat a cita agendada en menos de 24 horas.
\end{itemize}

\subsection{P2: alto impacto / esfuerzo medio-alto}
\begin{itemize}
  \item Configurar eventos de calidad post-chat (lead calificado, cita agendada, cita asistida) para cierre de embudo.
  \item Probar campana complementaria de trafico hacia landing de conversion con UTM y seguimiento.
  \item Construir biblioteca creativa con evidencia social local (testimonios autorizados y piezas educativas).
\end{itemize}

\subsection{P3: experimental}
\begin{itemize}
  \item Probar audiencias similares de conversaciones de alta calidad.
  \item Ejecutar test de posicionamiento por propuesta de valor (especializacion vs cercania vs rapidez de atencion).
  \item Explorar contenido colaborativo con instituciones educativas locales para autoridad y demanda incremental.
\end{itemize}

\section{Fuentes y limitaciones}
\subsection{Fuentes internas}
\begin{itemize}
  \item \texttt{report\_metadata}
  \item \texttt{page\_overview}
  \item \texttt{funnel}
  \item \texttt{traffic\_quality}
  \item \texttt{timeseries\_daily}
  \item \texttt{geo}
  \item \texttt{demographics}
  \item \texttt{campaigns\_available}
\end{itemize}

\subsection{Fuentes competitivas publicas}
\begin{itemize}
  \item \url{https://psicologiaymente.com/psicologos/2069561/dra-marizol-jimenez-luzon}
  \item \url{https://www.psico.org/centro-52334}
  \item \url{https://aprendiendojuntos.ec/}
  \item \url{https://tecnologicosudamericano.edu.ec/bienestar-estudiantil/centro-psicoterapeutico-la-cabana/}
  \item \url{https://www.terappio.com/ec/Psychologists/services/child-therapy/ecuador/loja}
  \item \url{https://psicologiaymente.com/directorio/ec/rankings/mejores-psicologos-loja}
\end{itemize}

\subsection{Fuentes visuales}
No se incorporaron capturas en esta corrida automatizada. Se deja preparada la ruta para evidencia visual en futuras versiones:
\texttt{docs/reports/figures/centro-de-psicologia-integral-y-terapia-psicotelcon/} (\texttt{observado}).

\subsection{Limitaciones}
\begin{itemize}
  \item No se ejecuto \texttt{apify-market-research}; posible causa: ausencia de token/creditos o no disponibilidad operativa en esta corrida.
  \item No se dispone de CRM cerrado de ventas/citas para calcular impacto economico final.
  \item La evidencia competitiva corresponde a senales publicas web/directorios y puede no reflejar totalmente la operacion offline de cada competidor.
  \item Evidencia visual parcial: pendiente de captura automatizada o carga manual de screenshots con fecha/hora y URL fuente.
\end{itemize}

\end{document}
